\IfFileExists{stacks-project.cls}{%
\documentclass{stacks-project}
}{%
\documentclass{amsart}
}
\usepackage{amssymb}

% For dealing with references we use the comment environment
\usepackage{verbatim}
\usepackage{fancyvrb}
\usepackage{pmboxdraw}
\newenvironment{reference}{\comment}{\endcomment}
%\newenvironment{reference}{}{}
\newenvironment{slogan}{\comment}{\endcomment}
\newenvironment{history}{\comment}{\endcomment}

% For source figures retained as assets
\usepackage{graphicx}

% For commutative diagrams we use Xy-pic
\usepackage[all]{xy}

% We use 2cell for 2-commutative diagrams.
\xyoption{2cell}
\UseAllTwocells

% We use multicol for the list of chapters between chapters
\usepackage{multicol}

% This is generally recommended for better output
\usepackage{lmodern}
\usepackage[T1]{fontenc}
\usepackage{textalpha}

% For cross-file-references
% We prefer the snapshot tmp/xref/#2.aux,
% which latexmkrc replaces atomically, so
% that documents compiled simultaneously
% never read a half-written .aux file.
% The last argument keeps links to #2.pdf.
\usepackage{xr-hyper}
\newcommand{\maybeexternaldocument}[2]{%
  \IfFileExists{#2.aux}{%
    \externaldocument[#1]{#2}%
  }{%
    \PackageWarning{stack}{Missing #2.aux; external references with prefix #1 unavailable}%
  }%
}
\let\externaldocumentorig\externaldocument
\renewcommand{\externaldocument}[2][]{%
  \IfFileExists{tmp/xref/#2.aux}{%
    \externaldocumentorig[#1]{tmp/xref/#2}[#2.pdf]%
  }{%
  \IfFileExists{#2.aux}{%
    \externaldocumentorig[#1]{#2}%
  }{%
    \PackageWarning{stack}{Missing #2.aux; external references with prefix #1 unavailable}%
  }%
  }%
}

% Package for hypertext links:
\usepackage{hyperref}

% For any local file, say "hello.tex" you want to link to please
% use \externaldocument[hello-]{hello}
\externaldocument[introduction-]{introduction}
\externaldocument[conventions-]{conventions}
\externaldocument[sets-]{sets}
\externaldocument[categories-]{categories}
\externaldocument[topology-]{topology}
\externaldocument[sheaves-]{sheaves}
\externaldocument[sites-]{sites}
\externaldocument[stacks-]{stacks}
\externaldocument[fields-]{fields}
\externaldocument[algebra-]{algebra}
\externaldocument[brauer-]{brauer}
\externaldocument[homology-]{homology}
\externaldocument[derived-]{derived}
\externaldocument[simplicial-]{simplicial}
\externaldocument[more-algebra-]{more-algebra}
\externaldocument[smoothing-]{smoothing}
\externaldocument[modules-]{modules}
\externaldocument[sites-modules-]{sites-modules}
\externaldocument[injectives-]{injectives}
\externaldocument[cohomology-]{cohomology}
\externaldocument[sites-cohomology-]{sites-cohomology}
\externaldocument[dga-]{dga}
\externaldocument[dpa-]{dpa}
\externaldocument[sdga-]{sdga}
\externaldocument[hypercovering-]{hypercovering}
\externaldocument[schemes-]{schemes}
\externaldocument[constructions-]{constructions}
\externaldocument[properties-]{properties}
\externaldocument[morphisms-]{morphisms}
\externaldocument[coherent-]{coherent}
\externaldocument[divisors-]{divisors}
\externaldocument[limits-]{limits}
\externaldocument[varieties-]{varieties}
\externaldocument[topologies-]{topologies}
\externaldocument[descent-]{descent}
\externaldocument[perfect-]{perfect}
\externaldocument[more-morphisms-]{more-morphisms}
\externaldocument[flat-]{flat}
\externaldocument[groupoids-]{groupoids}
\externaldocument[more-groupoids-]{more-groupoids}
\externaldocument[etale-]{etale}
\externaldocument[chow-]{chow}
\externaldocument[intersection-]{intersection}
\externaldocument[pic-]{pic}
\externaldocument[weil-]{weil}
\externaldocument[adequate-]{adequate}
\externaldocument[dualizing-]{dualizing}
\externaldocument[duality-]{duality}
\externaldocument[discriminant-]{discriminant}
\externaldocument[derham-]{derham}
\externaldocument[local-cohomology-]{local-cohomology}
\externaldocument[algebraization-]{algebraization}
\externaldocument[curves-]{curves}
\externaldocument[resolve-]{resolve}
\externaldocument[models-]{models}
\externaldocument[functors-]{functors}
\externaldocument[equiv-]{equiv}
\externaldocument[pione-]{pione}
\externaldocument[etale-cohomology-]{etale-cohomology}
\externaldocument[proetale-]{proetale}
\externaldocument[relative-cycles-]{relative-cycles}
\externaldocument[more-etale-]{more-etale}
\externaldocument[trace-]{trace}
\externaldocument[crystalline-]{crystalline}
\externaldocument[spaces-]{spaces}
\externaldocument[spaces-properties-]{spaces-properties}
\externaldocument[spaces-morphisms-]{spaces-morphisms}
\externaldocument[decent-spaces-]{decent-spaces}
\externaldocument[spaces-cohomology-]{spaces-cohomology}
\externaldocument[spaces-limits-]{spaces-limits}
\externaldocument[spaces-divisors-]{spaces-divisors}
\externaldocument[spaces-over-fields-]{spaces-over-fields}
\externaldocument[spaces-topologies-]{spaces-topologies}
\externaldocument[spaces-descent-]{spaces-descent}
\externaldocument[spaces-perfect-]{spaces-perfect}
\externaldocument[spaces-more-morphisms-]{spaces-more-morphisms}
\externaldocument[spaces-flat-]{spaces-flat}
\externaldocument[spaces-groupoids-]{spaces-groupoids}
\externaldocument[spaces-more-groupoids-]{spaces-more-groupoids}
\externaldocument[bootstrap-]{bootstrap}
\externaldocument[spaces-pushouts-]{spaces-pushouts}
\externaldocument[spaces-chow-]{spaces-chow}
\externaldocument[groupoids-quotients-]{groupoids-quotients}
\externaldocument[spaces-more-cohomology-]{spaces-more-cohomology}
\externaldocument[spaces-simplicial-]{spaces-simplicial}
\externaldocument[spaces-duality-]{spaces-duality}
\externaldocument[formal-spaces-]{formal-spaces}
\externaldocument[restricted-]{restricted}
\externaldocument[spaces-resolve-]{spaces-resolve}
\externaldocument[formal-defos-]{formal-defos}
\externaldocument[defos-]{defos}
\externaldocument[cotangent-]{cotangent}
\externaldocument[examples-defos-]{examples-defos}
\externaldocument[algebraic-]{algebraic}
\externaldocument[examples-stacks-]{examples-stacks}
\externaldocument[stacks-sheaves-]{stacks-sheaves}
\externaldocument[criteria-]{criteria}
\externaldocument[artin-]{artin}
\externaldocument[quot-]{quot}
\externaldocument[stacks-properties-]{stacks-properties}
\externaldocument[stacks-morphisms-]{stacks-morphisms}
\externaldocument[stacks-limits-]{stacks-limits}
\externaldocument[stacks-cohomology-]{stacks-cohomology}
\externaldocument[stacks-perfect-]{stacks-perfect}
\externaldocument[stacks-introduction-]{stacks-introduction}
\externaldocument[stacks-more-morphisms-]{stacks-more-morphisms}
\externaldocument[stacks-geometry-]{stacks-geometry}
\externaldocument[moduli-]{moduli}
\externaldocument[moduli-curves-]{moduli-curves}
\externaldocument[examples-]{examples}
\externaldocument[exercises-]{exercises}
\externaldocument[guide-]{guide}
\externaldocument[desirables-]{desirables}
\externaldocument[coding-]{coding}
\externaldocument[obsolete-]{obsolete}
\externaldocument[fdl-]{fdl}
\externaldocument[index-]{index}

% Heap Project documents
\externaldocument[k3-]{k3}
\externaldocument[torelli-]{torelli}
\externaldocument[linear-algebra-]{linear-algebra}
\externaldocument[groups-]{groups}
\externaldocument[complex-analysis-]{complex-analysis}
\externaldocument[functional-analysis-]{functional-analysis}
\externaldocument[categories-]{categories}
\externaldocument[commutative-algebra-]{commutative-algebra}
\externaldocument[ringed-spaces-]{ringed-spaces}
\externaldocument[homological-algebra-]{homological-algebra}
\externaldocument[lie-algebras-]{lie-algebras}
\externaldocument[representation-theory-]{representation-theory}
\externaldocument[smooth-manifolds-]{smooth-manifolds}
\externaldocument[differential-topology-]{differential-topology}
\externaldocument[algebraic-topology-]{algebraic-topology}
\externaldocument[differential-geometry-]{differential-geometry}
\externaldocument[algebraic-geometry-]{algebraic-geometry}
\externaldocument[symplectic-geometry-]{symplectic-geometry}
\externaldocument[complex-geometry-]{complex-geometry}
\externaldocument[hodge-theory-]{hodge-theory}
\externaldocument[%
variations-of-hodge-structures-]{%
variations-of-hodge-structures}
\externaldocument[%
homological-mirror-symmetry-]{%
homological-mirror-symmetry}
\maybeexternaldocument{deformations-}{deformations}
\externaldocument[vertex-algebras-]{vertex-algebras}
\externaldocument[springer-]{springer}
\externaldocument[infty-categories-]{infty-categories}
\externaldocument[seminars-differential-geometry-]{%
seminars-differential-geometry}
\externaldocument[seminars-complex-geometry-]{%
seminars-complex-geometry}
\externaldocument[seminars-algebraic-geometry-]{%
seminars-algebraic-geometry}
\externaldocument[seminars-analysis-pde-]{%
seminars-analysis-pde}
\externaldocument[surfaces-]{surfaces}
\externaldocument[birational-maps-commutative-algebra-]{birational-maps-commutative-algebra}
\externaldocument[cremona-transformations-]{cremona-transformations}
\externaldocument[derived-categories-of-sheaves-]{derived-categories-of-sheaves}
\externaldocument[geometric-invariant-theory-]{geometric-invariant-theory}
\externaldocument[geometric-stability-conditions-and-group-actions-]{geometric-stability-conditions-and-group-actions}
\externaldocument[moduli-spaces-of-sheaves-]{moduli-spaces-of-sheaves}
\externaldocument[bridgeland-stability-general-theory-]{bridgeland-stability-general-theory}
\externaldocument[hilbert-schemes-]{hilbert-schemes}
\externaldocument[noncommutative-abelian-surfaces-and-kummer-type-hyperkahler-varieties-]{noncommutative-abelian-surfaces-and-kummer-type-hyperkahler-varieties}
\externaldocument[mumford-tate-groups-in-hodge-theory-]{mumford-tate-groups-in-hodge-theory}
\externaldocument[hodge-birational-atoms-2026-]{hodge-birational-atoms-2026}
\externaldocument[xxii-egd-2026-]{%
xxii-egd-2026}
\externaldocument[pde-]{pde}
\externaldocument[probability-]{probability}

\externaldocument[linguistics-]{linguistics}

\externaldocument[%
16th-alga-meeting-2026-]{%
16th-alga-meeting-2026}

\externaldocument[%
iv-jornadas-lefschetz-2026-]{%
iv-jornadas-lefschetz-2026}

\externaldocument[reading-the-masters-]{%
reading-the-masters}

\externaldocument[real-analysis-]{%
real-analysis}

\externaldocument[optimal-transport-]{%
optimal-transport}

\externaldocument[history-]{history}

\externaldocument[seiberg-witten-]{seiberg-witten}

\externaldocument[history-philosophy-biology-]{history-philosophy-biology}

\externaldocument[classical-mechanics-]{classical-mechanics}
\externaldocument[electromagnetism-]{electromagnetism}
\externaldocument[waves-and-optics-]{waves-and-optics}
\externaldocument[wif-]{waves-in-fluids-iv}
\externaldocument[thermal-physics-]{thermal-physics}
\externaldocument[quantum-mechanics-]{quantum-mechanics}
\externaldocument[special-relativity-]{special-relativity}
\externaldocument[physics-]{physics}

\externaldocument[chemistry-]{chemistry}

\externaldocument[molecular-cell-biology-]{molecular-cell-biology}

\externaldocument[microbiology-]{microbiology}

\externaldocument[protistology-]{protistology}

\externaldocument[mycology-]{mycology}

\externaldocument[botany-]{botany}

\externaldocument[zoology-]{zoology}

\externaldocument[animal-histology-]{animal-histology}

\externaldocument[animal-physiology-]{animal-physiology}

\externaldocument[systematics-biogeography-]{systematics-biogeography}

\externaldocument[ecology-natural-resources-]{ecology-natural-resources}

\externaldocument[earth-sciences-]{earth-sciences}

\externaldocument[biostatistics-]{biostatistics}

\externaldocument[mathematical-systems-biology-]{mathematical-systems-biology}

\externaldocument[cimpa-el-salvador-2024-]{cimpa-el-salvador-2024}

\externaldocument[anthropology-]{anthropology}

\externaldocument[yang-mills-theory-]{yang-mills-theory}

\externaldocument[seminars-algebra-]{seminars-algebra}

\externaldocument[seminars-physics-]{seminars-physics}

\externaldocument[seminars-computing-]{seminars-computing}

\externaldocument[oktobermat-2026-]{oktobermat-2026}

\externaldocument[saet-2026-]{saet-2026}

\externaldocument[seminars-symplectic-geometry-]{seminars-symplectic-geometry}

\externaldocument[seminars-dynamics-]{seminars-dynamics}

\externaldocument[oktobermat-2024-]{oktobermat-2024}

\externaldocument[quantum-groups-]{quantum-groups}

\externaldocument[legal-2027-]{legal-2027}

\externaldocument[derived-categories-]{derived-categories}

% Theorem environments.
%
\theoremstyle{plain}
\newtheorem{theorem}[subsection]{Theorem}
\newtheorem{proposition}[subsection]{Proposition}
\newtheorem{lemma}[subsection]{Lemma}
\newtheorem{corollary}[subsection]{Corollary}

\theoremstyle{definition}
\newtheorem{definition}[subsection]{Definition}
\newtheorem{example}[subsection]{Example}
\newtheorem{exercise}[subsection]{Exercise}
\newtheorem{situation}[subsection]{Situation}

\theoremstyle{remark}
\newtheorem{remark}[subsection]{Remark}
\newtheorem{remarks}[subsection]{Remarks}

\numberwithin{equation}{subsection}

% Macros
%
\def\lim{\mathop{\mathrm{lim}}\nolimits}
\def\colim{\mathop{\mathrm{colim}}\nolimits}
\def\Spec{\mathop{\mathrm{Spec}}}
\def\Hom{\mathop{\mathrm{Hom}}\nolimits}
\def\Ext{\mathop{\mathrm{Ext}}\nolimits}
\def\SheafHom{\mathop{\mathcal{H}\!\mathit{om}}\nolimits}
\def\SheafExt{\mathop{\mathcal{E}\!\mathit{xt}}\nolimits}
\def\Sch{\mathit{Sch}}
\def\Mor{\mathop{\mathrm{Mor}}\nolimits}
\def\Ob{\mathop{\mathrm{Ob}}\nolimits}
\def\Sh{\mathop{\mathit{Sh}}\nolimits}
\def\NL{\mathop{N\!L}\nolimits}
\def\CH{\mathop{\mathrm{CH}}\nolimits}
\def\proetale{{pro\text{-}\acute{e}tale}}
\def\etale{{\acute{e}tale}}
\def\QCoh{\mathit{QCoh}}
\def\Ker{\mathop{\mathrm{Ker}}}
\def\Im{\mathop{\mathrm{Im}}}
\def\Coker{\mathop{\mathrm{Coker}}}
\def\Coim{\mathop{\mathrm{Coim}}}

% Boxtimes
%
\DeclareMathSymbol{\boxtimes}{\mathbin}{AMSa}{"02}

%
% Macros for moduli stacks/spaces
%
\def\QCohstack{\mathcal{QC}\!\mathit{oh}}
\def\Cohstack{\mathcal{C}\!\mathit{oh}}
\def\Spacesstack{\mathcal{S}\!\mathit{paces}}
\def\Quotfunctor{\mathrm{Quot}}
\def\Hilbfunctor{\mathrm{Hilb}}
\def\Curvesstack{\mathcal{C}\!\mathit{urves}}
\def\Polarizedstack{\mathcal{P}\!\mathit{olarized}}
\def\Complexesstack{\mathcal{C}\!\mathit{omplexes}}
% \Pic is the operator that assigns to X its picard group, usage \Pic(X)
% \Picardstack_{X/B} denotes the Picard stack of X over B
% \Picardfunctor_{X/B} denotes the Picard functor of X over B
\def\Pic{\mathop{\mathrm{Pic}}\nolimits}
\def\Picardstack{\mathcal{P}\!\mathit{ic}}
\def\Picardfunctor{\mathrm{Pic}}
\def\Deformationcategory{\mathcal{D}\!\mathit{ef}}


\begin{document}

\title{Religion}
\maketitle

\phantomsection
\label{section-phantom}

\tableofcontents

\section{Scope and sources}
\label{religion-section-religion-method}

These notes separate four layers:
direct observation; museum labels and
curatorial texts; published scholarship;
interpretation and open questions.

There are three sections: Christianity
(Section~\ref{section-christianity}),
Candombl\'e
(Section~\ref{section-candomble}), and
Johrei (Section~\ref{section-johrei}).
Describing a belief does not assert its
truth or represent every practitioner's
understanding.

\section{Christianity}
\label{section-christianity}

\subsection{Rosary and Marian titles}
\label{religion-section-rosary-marian-titles}

The Rosary is both a Catholic form of
prayer and the string of beads used to
count its repeated prayers.  Our Lady of
the Rosary is Mary under a devotional
title.  She is not a different saint.

The same distinction applies to Our Lady
of Carmel, of the Conception, of Sorrows,
of the Good Death, and other invocations.
They are titles of Mary associated with
different doctrines, places, orders,
images, and devotional histories.

Our Lady of Carmel belongs to the
Carmelite tradition, whose name derives
from Mount Carmel.  The scapular is an
important Carmelite devotional object.
This tradition also explains Carmo
place-names in Salvador.

The Immaculate Conception concerns Mary's
own conception.  Catholic doctrine holds
that she was preserved from original sin
from its first instant.  This must not be
confused with the virginal conception of
Jesus.

\subsection{Catholic brotherhoods in Salvador}
\label{religion-section-rosario-dos-pretos}

The historical importance of the Church
of Our Lady of the Rosary of the Black
People lies in the Black Catholic
brotherhood associated with it.  The
brotherhood was erected in 1685, and
construction of the church began in the
early eighteenth century
\cite{ipac-rosario}.

A brotherhood was not merely a prayer
group.  Such institutions organized
worship, funerals, mutual aid, festivals,
and forms of collective belonging within
a slave society.

This raises a precise question: how did
enslaved and free Black people create
institutions of their own inside the
Catholic order of colonial Brazil?

\subsection{Jesuits and the Cathedral}
\label{religion-section-jesuits-cathedral}

The Jesuits are members of the Society of
Jesus, a Catholic religious order founded
by Ignatius of Loyola and his companions
and approved in 1540.

Education became one of their central
activities.  This makes the connection
between the colonial Jesuit college in
Salvador and modern Jesuit institutions
such as PUC--Rio and the Universidad
Iberoamericana in Mexico conceptually
real, although the institutions belong to
very different historical periods.

Jesuits arrived in Bahia with Tom\'e de
Sousa in 1549.  The present Cathedral
Basilica was built as the church of the
Jesuit college and later became the
cathedral.

A cathedral is the church that contains
the bishop's \emph{cathedra}, or seat.  The
word names an ecclesiastical function,
not simply a large or ornate building.

\subsection{Saints and archangels}
\label{%
religion-section-saints-archangels-dates}
\label{religion-section-saint-george}

Saint Michael is celebrated on September
29 together with the archangels Gabriel
and Raphael in the Roman Catholic
calendar.

Saint George is a Christian warrior saint.
The plant called
\emph{lan\c{c}a-de-S\~ao-Jorge} takes its
popular name from the lance-like form and
the saint's martial image.

The plant is also associated in popular
practice with protection.  Associations
between Saint George and Ogum or Ox\'ossi
vary by region and religious house.  They
must not be stated as universal equations.

This association is discussed from both
sides: here from the side of the saint,
and from the side of Candombl\'e in the
debate on syncretism of
Subsection~\ref{%
religion-section-manifesto-1983}.

\section{Candomblé}
\label{section-candomble}

\subsection{Candomblé in Bahia}
\label{religion-section-what-candomble-is}

Candombl\'e is a Brazilian religion.  It
formed in Brazil through the preservation,
reorganization, and transformation of
several African religious traditions under
slavery and diaspora.

This does not mean that its divine beings
or all its concepts originated in Brazil.
Orix\'as, voduns, ancestors, divination,
ritual languages, and many other elements
have African histories.  The religious
institutions called Candombl\'e are
Afro-Brazilian formations.

Three identifications must therefore be
avoided:

\begin{itemize}
\item Africa is not one culture;
\item Yoruba religion is not identical
with Candombl\'e;
\item Candombl\'e is not identical with
Umbanda.
\end{itemize}

Umbanda is also Brazilian, but it formed
through a later and different history,
strongly associated with urban southeastern
Brazil in the twentieth century
\cite{prandi1990umbanda}.

\subsection{Candomblé nations}
\label{religion-section-candomble-nations}

A Candombl\'e ``nation'' is not a modern
nation-state.  It is a historical-ritual
classification involving language,
lineage, liturgy, divine beings, music,
and practice.

A first map is:

\begin{itemize}
\item Ketu or Nag\^o traditions are mainly
connected with Yoruba inheritances and
orix\'as;
\item Jeje traditions are mainly connected
with Ewe--Fon inheritances and voduns;
\item Angola and Congo traditions are
mainly connected with Central African,
Bantu-language inheritances and
\emph{nkisi}, called inquices in Brazil.
\end{itemize}

This is only a first map.  Historical
labels include Ketu, Jeje, Nag\^o, Ijex\'a,
Congo, Angola, and others, and ritual
languages are one important marker among
several \cite{castro1981lingua}.

The modern country Angola and the
Candombl\'e nation Angola are therefore
related historically but are not the same
kind of object.

\subsection{Yoruba, Nagô, and Ketu}
\label{religion-section-yoruba-nago-ketu}

Yoruba names an African language and a
broad historical and cultural field.
Nag\^o is a diasporic category that became
especially important in Bahia.  It cannot
be treated as a mathematically exact
synonym for Yoruba.

Ketu is one major tradition inside the
broader Nag\^o/Yoruba field of Brazilian
Candombl\'e.  The prominence of Nag\^o
traditions in Bahia was itself historically
constructed rather than timeless
\cite{pares2010nago}.

No single orix\'a is simply ``the orix\'a
of Ketu''.  A Ketu house cultivates many
orix\'as.  Oxal\'a is a major orix\'a
associated with creation, while a
particular house can have a patron orix\'a.

\subsection{Supreme divinity and orixás}
\label{%
religion-section-supreme-divinity-orixas}

Ol\'od\`umar\`e and Olorum are names or
titles used for the supreme divinity in
Yoruba-derived cosmologies.  The orix\'as
are not simply interchangeable names for
this supreme divinity.

One should distinguish:

\begin{itemize}
\item \emph{Orun}, the spiritual or
invisible realm;
\item \emph{Aye} or \emph{Aiye}, the
visible earthly realm;
\item Olorum, the supreme divinity;
\item the orix\'as, such as Oxal\'a,
Xang\^o, Oxum, Ox\'ossi, Ogum, and others.
\end{itemize}

The film title
\emph{Orun Aiye---A cria\c{c}\~ao do mundo}
uses the first two terms.  A creation
narrative centred on Obatal\'a or Oxal\'a
belongs to one important body of
Yoruba religious narrative
\cite{marins2012obatala}.

\subsection{Orun, Aiye, and Orunmila}
\label{religion-section-orun-aiye-film}

The stop-motion film
\emph{Orun Aiye---A cria\c{c}\~ao do mundo}
presents this cosmological vocabulary on
screen.  It was directed by Jamile Coelho
and Cintia Maria.  In the voice cast,
Carlinhos Brown plays Oxal\'a, Jo\~ao
Miguel plays
Ol\'od\`umar\`e,
and Jorge Washington plays Orunmila
\cite{orun-aiye-sepromi}.

Jorge Washington is also a long-standing
actor of the Bando de Teatro Olodum.  This
places the film inside a wider Black
performing-arts history in Salvador, rather
than treating the names in the credits as
isolated celebrity appearances
\cite{bando-jorge-washington}.

\subsection{Candomblé and Santería}
\label{religion-section-candomble-santeria}

Cuban Santer\'\i a, or Regla de Ocha, also
cultivates divine beings called
\emph{orichas}.  This does not make it the
same religion as Candombl\'e.

The useful comparison is historical:
Yoruba religious inheritances were
reorganized in different American
societies.  Nag\^o in Brazil and Lucum\'\i
in Cuba are diasporic categories with
different histories
\cite{lovejoyojo2015lucumi}.

\subsection{Terreiro and leadership}
\label{religion-section-terreiro-leadership}

A terreiro is not merely a building.  It is
a sacred territory, a ritual institution,
a community, and a network of initiated
relations.

An \emph{ialorix\'a} is the female
sacerdotal leader of a terreiro.  The
masculine counterpart is a
\emph{babalorix\'a}.  An ialorix\'a leads a
house; she does not command an entire
Candombl\'e nation.

There is no single central authority that
governs all Candombl\'e houses.  Houses can
be joined by initiatory genealogies,
alliances, influence, and shared ritual
histories while retaining their own
leadership and singularity.

\subsection{Ilê Axé Opô Afonjá}
\label{religion-section-opo-afonja}

Il\^e Ax\'e Op\^o Afonj\'a is a major
Ketu/Nag\^o terreiro in Salvador.  M\~ae
Aninha founded it in 1910 in a lineage
connected with Casa Branca do Engenho
Velho.  Xang\^o Afonj\'a is the patron of
the house \cite{barata2025ancestralidade}.

Its current leadership belongs to M\~ae Ana
de Xang\^o, the sixth ialorix\'a in the
line presented by the 2025 book.

The book
\emph{Pedagogia da ancestralidade}
should be read as a deep study of one
especially influential house, not as a
universal manual of Candombl\'e
\cite{barata2025ancestralidade}.

Its central idea is a pedagogy of the
terreiro.  Knowledge is transmitted through
coexistence, observation, participation,
oral memory, hierarchy, practice, care,
and the example of more experienced
members.

The house also developed civil,
educational, and cultural institutions,
including a school, library, Casa do Alak\'a,
and Museu Il\^e Ohun Lailai.  These do not
turn the terreiro into a secular museum;
they show that ritual, education, memory,
and social organization coexist.

\subsection{Mystery and limits of access}
\label{religion-section-mystery-limits}

The Afonj\'a book insists on preserving the
mystery of ritual life.  Photographs and
recordings may be authorized in community
contexts while remaining forbidden during
sacred rites
\cite{barata2025ancestralidade}.

This matters for visitors.  A
terreiro is not a tourist attraction whose
meaning is exhausted by access.  Permission,
conduct, photography, and the difference
between public and restricted spaces are
part of the religious reality.

\subsection{Mãe Stella de Oxóssi}
\label{religion-section-mae-stella}

M\~ae Stella de Ox\'ossi, Maria Stella de
Azevedo Santos, lived from 1925 to 2018.
She led Op\^o Afonj\'a as its fifth
ialorix\'a from 1976.

Three timeline panels at MUNCAB set out
her public and intellectual trajectory:

\begin{itemize}
\item she was initiated by M\~ae Senhora in
1939 and received the name Od\'e Kayod\^e;
\item she worked for decades in nursing and
public health;
\item in 1981 she travelled to Yoruba
centres in Nigeria, including Oshogbo,
Il\'e-If\`e, and Ed\'e;
\item in 1983 she intervened publicly in
the debate on syncretism during the
international orix\'a conference in
Salvador;
\item in 1986 she participated in the next
world conference in New York;
\item in 1987 she travelled to Benin in a
Brazilian cultural delegation organized by
Pierre Verger;
\item \emph{Meu tempo \'e agora} was
published in the early 1990s; Bahian
institutional records date the first
edition to 1993;
\item in 1999 the terreiro was protected by
IPHAN during her leadership;
\item biographical sources record that she
was among the Brazilian women indicated in
the 2005 Nobel Peace Prize initiative;
\item in 2013 she became the first
ialorix\'a and first Black woman to occupy
a chair in the Academia de Letras da Bahia;
\item her death in 2018 was described in
religious language as a return to Orun.
\end{itemize}

These milestones match the MUNCAB timeline
and Bahian institutional biographical
records \cite{alba-mae-stella,mae-stella-alb}.
They explain why the Op\^o Afonj\'a book is
not merely a local curiosity: it gives
access to the house in which one of the
most publicly and intellectually prominent
ialorix\'as of recent Brazilian history
worked and led.

It is still better to say ``one of the most
prominent'' than to declare M\~ae Stella
the single most famous ialorix\'a, since
that would require an artificial ranking.

\subsection{The 1983 manifesto}
\label{religion-section-manifesto-1983}

MUNCAB reproduced the manifesto
\emph{O p\'ublico e o povo do Candombl\'e},
published in Salvador on August 12, 1983.
The panel attributed it to Stella de
Ox\'ossi, Menininha do Gantois, Tet\^e de
Ians\~a, Olga de Alaketu, and Nicinha do
Bogum.

The manifesto is a primary source.  It
rejects descriptions of Candombl\'e as a
sect or primitive animism.  It also rejects
the idea that Candombl\'e requires Catholic
forms in order to be legitimate.

Its argument about syncretism is historical
and political.  It connects concealment
and Catholic forms with slavery and
submission, while demanding the right to
practise Candombl\'e in its own terms.

It also attacks folklorization,
commercialization, press sensationalism,
and tourist consumption.  Its reference to
the Lavagem do Bonfim shows that a practice
can be religiously and politically debated
by Candombl\'e leaders themselves.

The manifesto does not claim that
syncretism never existed.  It disputes what
that history should require in the present.

\subsection{Water and maternity in Omi Tutu}
\label{religion-section-omi-tutu-religion}

Hariel Revignet's MUNCAB exhibition
\emph{Omi Tutu---\'Agua fresca} placed
water in relation with territory, body,
displacement, continuity, and return
\cite{muncab2026omitutu}.  The curatorial
text also brought Oxum, Iemanj\'a, and
Nan\~a into a field of fresh water, sea,
maternity, memory, gestation, and
transformation.

This is useful religious vocabulary, but
a curatorial text is not a complete
theology of the orix\'as.  The presence of
an orix\'a in an artwork also does not
make the work specifically Umbandist:
orix\'as belong to Yoruba-derived
religious worlds and are central to
Candombl\'e traditions as well.  The
exhibition presents those names and
associations through an artistic
interpretation, which should be compared
with sources internal to the relevant
Candombl\'e traditions.

\subsection{Ofá and Oxóssi}
\label{religion-section-ofa-oxossi}

A small iron sculpture by Diogum
represents the \emph{of\'a}, the bow and
arrow associated with Ox\'ossi, by its
curved form and central shaft.

The heading ``Di Ogum'' names the artist,
Diogo Oliveira, and his relation to Ogum.
It does not mean that every symbol in his
work represents Ogum.

The case is useful because material,
artist's identity, and represented attribute
can point in different directions.

\subsection{Airá and Oxalufã}
\label{religion-section-aira-oxalufa}

Airá (also written Ayrá) is not a form
of Ogum. He is associated with Xangô:
the name ``Xangô Airá'' is found, but
some traditions recognize him as a
distinct orixá.\footnote{
Correio Nagô, ``Ayrá: o manto branco
da justiça'', 29 June 2015:
\href{https://correionago.com.br/%
ayra-o-manto-branco-da-justica/}
{article}; and the web page cited below.
Accessed 15 October 2026.}

According to the web page of one
terreiro, in one version of the story
Oxalufã, the old Oxalá, spends seven
years unjustly imprisoned in the kingdom
of Xangô, without Xangô knowing. Xangô
discovers the injustice and frees him.
As Oxalufã is weakened, Xangô asks Airá
to accompany him on the way home. Airá
carries him on his back, lights fires to
warm him, and tells him stories. They
become companions, in a relation
described as that of father and
son.\footnote{
Ilê Àsé Oni Èwé Ossayn Atin L'Oyá,
``Ayrá (Airá)'', section
``História de Ayrá'':
\href{https://www.bruxodailha.com.br/%
orixas/ayra/}{web page}.
Accessed 15 October 2026.}

In this version, Xangô frees; Airá
carries and cares. The details vary
between narratives and terreiros.
According to the same page, Airá has a
key linked to the weather: this is
recorded as an attribution of that
source, not as a universal symbol nor as
a safe identification of an isolated
image.

\subsection{Beads and cowries}
\label{religion-section-beads-cowries}

A bead is first a material object.  Its
religious meaning comes from arrangement,
colour, preparation, use, lineage, and
context.

Strings of beads can identify religious
relations and may be ritually prepared.
One cannot infer an orix\'a mechanically
from any necklace or colour combination.

Three works make the material central:
Hariel Revignet's bead installation,
N\'adia Taquary's strings in the colours of
Iroko, and Sandro Aiy\^e's wooden sculpture
with beads.

Cowries are actual materials in Hariel
Revignet's works.  Their visual and
religious histories require a separate
study; they should not be reduced to a
generic sign of ``Africa''.

\subsection{Iroko and Iyami}
\label{religion-section-iroko-iyami}

N\'adia Taquary's
\emph{Ir\^ok\^o: A \'arvore c\'osmica}
connected Iroko with time, ancestry, the
life cycle, the gameleira, and the Iyami.

The museum label renders Iyami as
``feiticeiras''.  That translation is too
narrow if taken as a definition.  The
concept belongs to a larger Yoruba field of
ancestral maternal and feminine power
\cite{santos2008iyami}.

The comparison with Circe is legitimate as
a comparative question about female power,
transformation, and fear.  It does not
establish transmission or historical
kinship between Greek and Yoruba myths.

\subsection{Bará, Exu, and paths}
\label{religion-section-bara-exu-paths}

Sandro Aiy\^e's sculpture
\emph{B\'ar\`a II} was shown in the MUNCAB
exhibition
\emph{Pad\^e On\~a---Encontrar Caminhos}.
The official framing presents Exu through
movement, mediation, circulation, and the
opening of paths
\cite{muncab2026padeona}.

This is better than identifying the work
only from a trident-like object.  An
attribute can call attention, but the
exhibition title and curatorial context are
stronger evidence.

\subsection{Bembé and public religion}
\label{%
religion-section-bembe-public-religion}

The book on the Bemb\'e do Mercado opens a
question distinct from the life of a single
terreiro: how does Candombl\'e organize
sacred territory and collective action in
public space?

The Bemb\'e is celebrated in Santo Amaro
and involves many terreiros, public ritual,
heritage, urban territory, and memories of
abolition \cite{ufrb2026bembe}.

It belongs to the study of religion, but
also to the histories of the Rec\^oncavo,
public space, and cultural heritage.

\subsection{Further reading}
\label{religion-section-religion-reading}

\begin{enumerate}
\item \emph{Pedagogia da ancestralidade},
a study of Op\^o Afonj\'a
\cite{barata2025ancestralidade};
\item a reliable panoramic account of
Candombl\'e nations;
\item \emph{Na\c{c}\~oes do Candombl\'e: Jeje},
a deliberately particular Jeje study.
The volume identifies itself as Book 1
of the larger
\emph{Na\c{c}\~oes do Candombl\'e}
project; it should not be mistaken for a
general survey of the whole religion
\cite{ufrb2026nacoes};
\item a comparison of Ketu/Nag\^o, Jeje,
and Angola/Congo that does not treat one
as the whole religion;
\item a historical comparison of
Candombl\'e with Umbanda and Santer\'\i a.
\end{enumerate}

\section{Johrei}
\label{section-johrei}

Johrei is a central practice of the
Igreja Messiânica Mundial do Brasil
(Church of World Messianity of Brazil).
The word is Japanese and is usually
translated as purification of the
spirit. In the Church's own description,
the person who ministers Johrei raises
the palm of the hand towards the person
who receives it, without touching, and
through it a divine light is transmitted
that purifies the spirit. Johrei is
ministered free of charge, in the units
of the Church and in homes.

The Church was founded in Japan by Mokiti
Okada (1882--1955), whom the members call
Meishu-Sama. The foundation is dated to
1 January 1935, under the name Dai Nippon
Kannon Kai; the present name dates from
1950. The Church arrived in Brazil in the
1950s, and its sacred place in the
country is the Solo Sagrado de
Guarapiranga, in São Paulo.

Johrei is taught together with two other
practices: Nature Farming
(\emph{Agricultura Natural}), which
cultivates without chemical fertilizers
or pesticides, and the Beautiful
(\emph{o Belo}), the cultivation of art
and beauty, of which ikebana is the
best-known example. The three are
presented as paths to what the doctrine
calls Paradise on Earth.

From the point of view of the study of
religions, the Church is one of the
Japanese new religions, and Johrei
belongs to the family of healing
practices without touch. The therapeutic
efficacy of Johrei is a statement of
faith; this note describes it and does
not assert it.

This summary is based on Brazilian
academic articles and has not been
checked against a printed source; the
description of the gesture follows the
Church's own account.

\bibliographystyle{amsalpha}
\bibliography{my}

\end{document}
